\subsection{柔软层}

定义 5. --- 设 \(p\colon E \rightarrow B\) 为平展映射。称映射 \(p\) 是柔软的，或称平展 B-空间 \((E,p)\) 是柔软的，如果 \(p\) 在 B 的闭子空间上的每个连续截面都能扩张为 \(p\) 在 B 上的一个连续截面。

设 F 为 B 上的一个层。若相伴的平展空间 (I, p. 50, def. 4) 是柔软的，则称 F 为柔软层。

设 \(\mathcal{F}\) 为 B 上的层。层 \(\mathcal{F}\) 是柔软的，当且仅当：对 B 的每个闭集 Z、Z 的每个开邻域 U 以及每个 \(s \in \mathcal{F}(\mathrm{U})\) ，存在 \(t \in \mathcal{F}(\mathrm{B})\) 以及 Z 的一个含于 U 的开邻域 V，使得 \(s \mid \mathrm{V} = t \mid \mathrm{V}\) 。

若 \(\mathcal{F}\) 是柔软层，则 \(\mathcal{F}(\mathrm{B})\) 非空：事实上，平展空间 \(\mathrm{E}_{\mathcal{F}}\) （与层 \(\mathcal{F}\) 相伴者）在 \(\varnothing\) 上的唯一截面可扩张为 \(\mathrm{E}_{\mathcal{F}}\) 在 B 上的连续截面。

设 \(p \colon \mathrm{E} \to \mathrm{B}\) 为平展映射，并设 A 为 B 的闭子空间。若 \(p\) 是柔软的，则映射 \(p_{\mathrm{A}} \colon \overline{p}^{1}(\mathrm{A}) \to \mathrm{A}\) 是柔软的。等价地，若 \(\mathcal{F}\) 是 B 上的柔软层，则它在闭子空间 A 上诱导的层是柔软的。

命题 6. --- 设 B 为拓扑空间，\(\mathcal{F}\) 为 B 上的层，\((\mathrm{A}_i)_{i\in \mathrm{I}}\) 为 B 的局部有限闭覆盖。层 \(\mathcal{F}\) 是柔软的，当且仅当对每个 \(i\in \mathrm{I}\)，诱导层 \(\mathcal{F}_{\mathrm{A}_i}\) 都是柔软的。

该条件显然是必要的。我们来证明它是充分的。用 \(p\colon \mathrm{E}\to \mathrm{B}\) 表示 B-平展空间 \(\mathrm{E}_{\mathcal{F}}\) （与层 \(\mathcal{F}\) 相伴者）。设 A 为 B 的闭子空间，并设 \(s\colon \mathrm{A}\to \mathrm{E}\) 为 \(p\) 在 A 上的一个连续截面；须证 \(s\) 具有到 B 上的连续扩张。对 I 的每个子集 J，令 \(\mathrm{A_J} = \bigcup_{i\in \mathrm{J}}\mathrm{A}_i\) ；集 \(\mathrm{A_J}\) 在 B 中是闭的 (TG, I, p. 6, prop. 4)。

设 \(\mathcal{S}\) 为由数对 \((\mathrm{J},t)\) 组成的集，其中 \(\mathrm{J}\) 是 I 的子集，而 \(t\) 是 E 在 \(\mathrm{A_J}\) 上的连续截面，且与 \(s\) 在 \(\mathrm{A} \cap \mathrm{A_J}\) 上重合。给 \(\mathcal{S}\) 赋以记作 \(\leqslant\) 的序关系：\((\mathrm{J},t) \leqslant (\mathrm{J}',t')\) 如果 \(\mathrm{J} \subset \mathrm{J}'\) 且 \(t' \mid \mathrm{A_J} = t\) 。对 \(\sigma = (\mathrm{J},t) \in \mathcal{S}\) ，记 \(\mathrm{J}_{\sigma} = \mathrm{J}\) 与 \(t_{\sigma} = t\) 。我们来证明序集 \(\mathcal{S}\) 是归纳的。设 S 为 \(\mathcal{S}\) 的一个全序子集。令 \(\mathrm{J} = \bigcup_{\sigma \in \mathrm{S}} \mathrm{J}_{\sigma}\) ；它是 I 的一个子集。于是便定义了截面 \(t\) （E 在 \(\mathrm{A_J}\) 上的），令 \(t(x) = t_{\sigma}(x)\) （当 \(x \in \mathrm{A}_{J_{\sigma}}\) 时）；因而有 \(t \mid \mathrm{A} \cap \mathrm{A_J} = s\) 。设 \(j \in \mathrm{J}\)，并设 \(\sigma \in S\) 使 \(j \in J_{\sigma}\) ；由于 \(t \mid A_j = t_{\sigma} \mid A_j\) ，\(t\) 在 \(A_j\) 上的限制是连续的。于是由 TG, I, p. 19, prop. 4 知 \(t\) 是连续的。这样，(J, t) 便是 \(\mathcal{S}\) 的元素；按构造，它是 S 的一个上界。这就证明了集 \(\mathcal{S}\) 是归纳的。因此它具有一个极大元 (J, t) (E, III, p. 20, th. 2)。

我们用反证法，假设 \(J \neq I\) 。设 i 为 \(I \setminus J\) 的一个元素。令 \(A' = (A_i \cap A) \cup (A_i \cap A_J)\)，并定义截面 \(s'\) （E 在 \(A'\) 上的）如下：

\[
s ^ {\prime} (a) = \left\{ \begin{array}{l l} s (a) & \text {for} a \in \mathrm{A} _ {i} \cap \mathrm{A}, \\ t (a) & \text {for} a \in \mathrm{A} _ {i} \cap \mathrm{A} _ {\mathrm{J}}, \end{array} \right.
\]

这是可能的，因为 \(s\) 与 \(t\) 在 \(\mathrm{A} \cap \mathrm{A}_{\mathrm{J}}\) 上重合。此外，由于 \(\mathrm{A}_{i} \cap \mathrm{A}\) 与 \(\mathrm{A}_{i} \cap \mathrm{A}_{\mathrm{J}}\) 都是闭集，截面 \(s'\) 是连续的 (TG, I, p. 19, prop. 4)。按假设，存在连续截面 \(s_{i}: \mathrm{A}_{i} \to \mathrm{E}\) 使之扩张 \(s'\) 。由于 \(s_{i}\) 与 \(t\) 在 \(\mathrm{A}_{\mathrm{J}} \cap \mathrm{A}_{i}\) 上的限制相等，映射 \(t': \mathrm{A}_{\mathrm{J} \cup \{i\}} \to \mathrm{E}\) （与 \(t\) 在 \(\mathrm{A}_{\mathrm{J}}\) 上重合、与 \(s_{i}\) 在 \(\mathrm{A}_{i}\) 上重合者）便是 \(p\) 在 \(\mathrm{A}_{\mathrm{J} \cup \{i\}}\) 上的连续截面，且它扩张 \(s| \mathrm{A} \cap \mathrm{A}_{\mathrm{J} \cup \{i\}}\) 。于是有 \((\mathrm{J}, t) < (\mathrm{J} \cup \{i\}, t')\) ，这与 \((\mathrm{J}, t)\) 为极大元的假设相矛盾。

于是 J = I，从而 \(A_{J} = B\)，且 t 是 E 在 B 上的连续截面，它扩张 s。

推论 1. --- 设 B 为仿紧空间，\(\mathcal{F}\) 为 B 上的层，\((\mathrm{U}_i)_{i\in \mathrm{I}}\) 为 B 的开覆盖。若对每个 \(i\in \mathrm{I}\) ，诱导层 \(\mathcal{F}|\mathrm{U}_i\) 都是柔软的，则层 \(\mathcal{F}\) 是柔软的。

事实上，存在局部有限闭覆盖 \((\mathrm{F}_{j})_{j\in\mathrm{J}}\) ，它比覆盖 \((\mathrm{U}_{i})_{i\in\mathrm{I}}\) 更细 (TG, IX, p. 49, prop. 4 及 p. 48, cor. 1)。于是，对每个 \(j\in J\)，层 \(\mathcal{F}|F_{j}\) 是柔软的，而命题蕴含层 \(\mathcal{F}\) 是柔软的。

推论 2. --- 设 B 为仿紧空间，\(\mathcal{F}\) 为 B 上的层，\((\mathrm{A}_i)_{i\in \mathrm{I}}\) 为 B 的局部有限闭覆盖。为使层 \(\mathcal{F}\) 是柔软的，必须且只需下述条件成立：

对每个 \(i \in I\) 、\(A_{i}\) 的每个闭子集 A、B 的每个含 A 的开集 V 以及 \(\mathcal{F}(V)\) 的每个元素 s，存在 \(A_{i}\) 在 B 中的一个开邻域 U、\(\mathcal{F}(U)\) 的一个元素 t 以及 A 在 B 中的一个含于 \(U \cap V\) 的开邻域 W，使得 \(t | W = s | W\) 。

设层 \(\mathcal{F}\) 是柔软的，我们来证明该条件成立。设 \(i\in\mathrm{I}\) ，设 A 为 \(\mathrm{A}_{i}\) 的闭子集，设 V 为 B 的含 A 的开集，并设 \(s\) 为 \(\mathcal{F}(\mathrm{V})\) 的元素。令 \(s_{0}=\sigma_{\mathcal{F}}(s)\) 。它是平展空间 \(\mathrm{E}_{\mathcal{F}}\) 在 V 上的连续截面。由于 \(\mathrm{A}_{i}\) 是闭的，A 在 B 中是闭的，故由柔软层的定义，\(s_{0}\mid\mathrm{A}\) 可扩张为一个截面 \(t_{0}\colon\mathrm{B}\to\mathrm{E}_{\mathcal{F}}\) 。截面 \(s_{0}\) 与 \(t_{0}\mid\mathrm{V}\) （V-平展空间 \(\mathrm{E}_{\mathcal{F}}\times_{\mathrm{B}}\mathrm{V}\) 的）在 A 上重合，从而在 A 的一个邻域 W 上重合 (I, p. 34, prop. 11, b))。

反之，设推论的条件成立。由命题 6，只需证明对每个 \(i \in I\)，层 \(\mathcal{F}_{A_i}\) 是柔软的。设 A 为 \(A_i\) 的闭子集，并设 \(s_0\) 为 \(\mathcal{F} | A_i\) 的平展空间在 A 上的截面，即平展空间 \(E_{\mathcal{F}}\) 在 A 上的连续截面。由于 A 在 B 中是闭的且 B 是仿紧的，\(s_0\) 可扩张为在 B 中 A 的一个邻域 V 上的截面 \(s\) (I, p. 37, th. 2)。按假设，存在 \(A_i\) 在 B 中的一个开邻域 U 以及一个截面 \(t\) （\(E_{\mathcal{F}}\) 在 U 上的），它在 A 的一个邻域上与 \(s\) 重合。\(t\) 在 \(A_i\) 上的限制是 \(\mathcal{F}_{A_i}\) 的一个截面，且它扩张 \(s_0\) 。于是层 \(\mathcal{F}_{A_i}\) 是柔软的。由命题 6，层 \(\mathcal{F}\) 是柔软的。

